\section*{Hoofdstuk \ref{ch:g11:cancer} --- Schade aan het genoom en kanker}

\begin{solution}{exo:g11:cancer:1}
\omterm{def:g11:cancer:cancer}{Kanker}: \omterm{def:g10:cells-common-unit:cell}{cellen} die zich zonder de normale regelingen delen en naburig
weefsel binnendringen. \omterm{def:g11:cancer:cancer}{Tumor}: de massa die zo'n kloon vormt. Uitzaaiing:
een nieuwe \omterm{def:g11:cancer:cancer}{tumor} die elders wordt gesticht door \omterm{def:g10:cells-common-unit:cell}{cellen} die via het bloed
of de lymfe zijn gereisd.
\end{solution}

\begin{solution}{exo:g11:cancer:2}
Een \omterm{def:g10:universal-dna:gene}{gen} waarvan het \omterm{def:g11:gene-expression:protein}{eiwit} groeisignalen doorgeeft en de \omterm{def:g10:cells-common-unit:cell}{cel} de cyclus in
duwt; is het zo gemuteerd dat het \omterm{def:g11:gene-expression:protein}{eiwit} blijvend werkzaam is, dan wordt
het een oncogen. Eén vastzittend gaspedaal drijft de \omterm{def:g10:cells-common-unit:cell}{cel} aan, wat het
andere \omterm{def:g10:universal-dna:gene}{allel} ook doet.
\end{solution}

\begin{solution}{exo:g11:cancer:3}
Het legt de cyclus stil wanneer het \omterm{def:g10:universal-dna:information}{DNA} beschadigd is en brengt het
herstel of de zelfvernietiging op gang. Eén werkend \omterm{def:g10:universal-dna:gene}{allel} maakt nog genoeg
\omterm{def:g11:gene-expression:protein}{eiwit} om te remmen; pas het verlies van beide neemt de rem weg.
\end{solution}

\begin{solution}{exo:g11:cancer:4}
Ultraviolet licht (huidkankers), teer uit tabak (longkanker), ioniserende
straling, asbest, alcohol, papillomavirus (baarmoederhalskanker).
\end{solution}

\begin{solution}{exo:g11:cancer:5}
Al haar \omterm{def:g10:cells-common-unit:cell}{cellen} stammen af van één \omterm{def:g10:cells-common-unit:cell}{cel} waarin de eerste \omterm{def:g11:mutations:mutation}{mutatie} optrad; zij
delen die \omterm{def:g11:mutations:mutation}{mutatie} en de \omterm{def:g11:mutations:mutation}{mutaties} die volgden, en kenmerken zoals hetzelfde
uitgeschakelde \omterm{def:g11:cell-cycle-mitosis:chromosome}{X-chromosoom}.
\end{solution}

\begin{solution}{exo:g11:cancer:6}
Ongeveer 200 op 40 en 2700 op 80: een factor 13 voor een factor 2 in
leeftijd --- verscheidene gebeurtenissen die zich vermenigvuldigen, geen
enkele evenredige oorzaak.
\end{solution}

\begin{solution}{exo:g11:cancer:7}
Ongeveer 10 en 17. Wie er 20 per dag rookt: $17 \times 1\% = 17\%$, één op
zes.
\end{solution}

\begin{solution}{exo:g11:cancer:8}
Zonder p53 houdt de \omterm{def:g10:cells-common-unit:cell}{cel} niet meer in om beschadigd \omterm{def:g10:universal-dna:information}{DNA} vóór de replicatie
te herstellen, en vernietigt zij zichzelf niet meer wanneer de schade
ernstig is: beschadigingen worden bij elke deling tot \omterm{def:g11:mutations:mutation}{mutaties} gekopieerd,
en \omterm{def:g10:cells-common-unit:cell}{cellen} met een overhoopgehaald genoom blijven leven.
\end{solution}

\begin{solution}{exo:g11:cancer:9}
Elke \omterm{def:g10:cells-common-unit:cell}{cel} van haar lichaam mist al één \omterm{def:g10:universal-dna:gene}{allel}; één enkele somatische \omterm{def:g11:mutations:mutation}{mutatie}
van het andere maakt het verlies compleet, in plaats van twee
onafhankelijke. Doordat er vanaf de geboorte één stap is gezet, worden de
overige stappen eerder en in meer \omterm{def:g10:cells-common-unit:cell}{cellen} bereikt. Haar zussen erven het
\omterm{def:g10:universal-dna:gene}{allel} van dezelfde ouder met kans $1/2$.
\end{solution}

\begin{solution}{exo:g11:cancer:10}
De middelen doden delende \omterm{def:g10:cells-common-unit:cell}{cellen} zonder onderscheid: de \omterm{def:g10:cells-common-unit:cell}{cellen} van de
haarwortels, de darmbekleding en het beenmerg delen zich snel en worden
samen met de \omterm{def:g11:cancer:cancer}{tumor} getroffen --- vandaar haaruitval, misselijkheid en een
tekort aan witte \omterm{def:g10:cells-common-unit:cell}{cellen}.
\end{solution}

\begin{solution}{exo:g11:cancer:11}
Stoppen beëindigt de aanvoer van nieuwe \omterm{def:g11:mutations:mutation}{mutaties}, en de \omterm{def:g10:cells-common-unit:cell}{cellen} van de
bekleding worden geleidelijk vervangen, waardoor veel klonen die vroege
stappen droegen verdwijnen. Maar sommige klonen met verscheidene stappen
blijven, en de \omterm{def:g11:mutations:mutation}{mutaties} die er al zijn worden niet uitgewist; het risico
blijft boven dat van een niet-roker.
\end{solution}

\begin{solution}{exo:g11:cancer:12}
$10^{12} \times 10^{-7} = 10^5$ \omterm{def:g10:cells-common-unit:cell}{cellen} lopen de eerste \omterm{def:g11:mutations:mutation}{mutatie} op. Elk
sticht een kloon die zich duizenden keren kan delen, dus wordt de tweede
\omterm{def:g11:mutations:mutation}{mutatie} niet in één \omterm{def:g10:cells-common-unit:cell}{cel} gezocht maar in alle delingen van de hele kloon,
en zo ook de derde: de stappen vermenigvuldigen het doelwit, en de
werkelijke kans ligt vele ordes van grootte boven $10^{-21}$.
\end{solution}

\begin{solution}{exo:g11:cancer:13}
Het viruseiwit neemt de rem weg zonder het \omterm{def:g10:universal-dna:gene}{gen} aan te raken: de besmette
\omterm{def:g10:cells-common-unit:cell}{cel} gedraagt zich alsof p53 verloren was en hoopt de \omterm{def:g11:mutations:mutation}{mutaties} op die het
proces voltooien. Het vaccin laat het afweersysteem het virus vernietigen
voordat het de baarmoederhals besmet; zijn de \omterm{def:g10:cells-common-unit:cell}{cellen} eenmaal besmet en
veranderd, dan helpt het vaccin niet meer --- vandaar inenting in de
vroege puberteit.
\end{solution}

\begin{solution}{exo:g11:cancer:14}
Een poliep is een kloon die één of twee stappen heeft gezet; hem
verwijderen neemt de \omterm{def:g10:cells-common-unit:cell}{cellen} weg waarin de overige stappen zouden kunnen
optreden, zodat de \omterm{def:g11:cancer:cancer}{kanker} nooit ontstaat. Omdat de reeks van poliep tot
\omterm{def:g11:cancer:cancer}{kanker} tien tot twintig jaar duurt, vangt een tussenpoos van tien jaar de
meeste poliepen op voordat zij haar voltooien.
\end{solution}

\begin{solution}{exo:g11:cancer:15}
Het voorkomen stijgt met de leeftijd omdat de stappen tientallen jaren
kosten, maar de stappen worden tijdens het leven gezet, grotendeels onder
blootstellingen --- alleen al tabak is goed voor een derde van de
sterfgevallen door \omterm{def:g11:cancer:cancer}{kanker}, en infecties, alcohol, zon en voeding voor een
groot deel van de rest. Een blootstelling wegnemen vermindert op elke
leeftijd de \omterm{def:g11:mutations:mutation}{mutaties} die nog moeten komen; screening haalt klonen weg
vóór de laatste stap. De ouderdom is wanneer de rekening wordt
gepresenteerd, niet wanneer zij wordt geschreven.
\end{solution}

\begin{solution}{pb:g11:cancer:1}
\textbf{1.} Zij zitten wel in de \omterm{def:g11:cancer:cancer}{tumor} en niet in de andere \omterm{def:g10:cells-common-unit:cell}{cellen} van de
patiënt, dus zijn zij in een lichaamscel ontstaan. Kiembaanmutaties zouden
in elke \omterm{def:g10:cells-common-unit:cell}{cel} zitten, het gezonde weefsel inbegrepen.

\textbf{2.} Dat de meeste \omterm{def:g11:mutations:mutation}{mutaties} van de \omterm{def:g11:cancer:cancer}{tumor} door benzopyreen zijn
veroorzaakt --- het \omterm{def:g11:mutations:mutagen}{mutageen} uit tabaksrook liet zijn handtekening achter.

\textbf{3.} Een \omterm{def:g11:mutations:mutation}{nonsense-mutatie}: een voortijdig stopcodon kapt p53 af,
dat daardoor niet kan werken.

\textbf{4.} p53 is een rem: één werkend \omterm{def:g10:universal-dna:gene}{allel} remt nog, dus moesten beide
verdwijnen (één gemuteerd, één verwijderd). De receptor is een gaspedaal:
één blijvend werkzaam \omterm{def:g11:gene-expression:protein}{eiwit} drijft de \omterm{def:g10:cells-common-unit:cell}{cel} aan, wat het andere \omterm{def:g10:universal-dna:gene}{allel} ook
doet.

\textbf{5.} Het receptorgen is het gaspedaal (een proto-oncogen dat
oncogen is geworden); p53 is de rem (tumorsuppressor).

\textbf{6.} $20 \times 365 \times 35 \approx \num{255000}$ sigaretten.

\textbf{7.} Ongeveer $\num{255000}$ substituties per \omterm{def:g10:cells-common-unit:cell}{cel} --- vijftig keer
de 5000 die zijn gevonden. De meeste beschadigingen worden hersteld, en
\omterm{def:g10:cells-common-unit:cell}{cellen} met te veel \omterm{def:g11:mutations:mutation}{mutaties} sterven; 5000 is het deel dat overbleef.

\textbf{8.} Zij traden vroeg op, in de stichtende \omterm{def:g10:cells-common-unit:cell}{cel} of haar eerste
nakomelingen, zodat elke tumorcel ze erfde; de andere ontstonden later in
deelklonen.

\textbf{9.} $10^{10} \approx 2^{33}$: 33 verdubbelingen, bij 100 dagen per
verdubbeling ongeveer 9 jaar.

\textbf{10.} Heeft de \omterm{def:g11:cancer:cancer}{tumor} 9 jaar of meer nodig gehad om uit zijn
stichtende \omterm{def:g10:cells-common-unit:cell}{cel} te groeien, dan trad de laatste sleutelmutatie ongeveer 25
jaar na het begin van de 35 rookjaren op, of eerder.

\textbf{11.} $17 \times 1\% = 17\%$.

\textbf{12.} 17 \omterm{def:g11:cancer:cancer}{kankers}, waarvan er 1 hoe dan ook zou zijn opgetreden: 16
op 100 zijn aan het roken toe te schrijven.

\textbf{13.} Er komen geen nieuwe \omterm{def:g11:mutations:mutation}{mutaties} bij; de bekleding vernieuwt
zichzelf en veel klonen die vroege stappen droegen, worden in de loop van
de jaren afgestoten en vervangen; de overblijvende klonen hebben minder
stappen gezet dan zij bij doorroken zouden hebben.

\textbf{14.} Verscheidene zeldzame \omterm{def:g11:mutations:mutation}{mutaties} moeten zich in één lijn
ophopen, en die lijn moet daarna tot een waarneembare omvang uitgroeien:
elk stadium kost jaren, en hun som bedraagt tientallen jaren.

\textbf{15.} Eén \omterm{def:g11:mutations:mutation}{mutatie} per \omterm{def:g10:cells-common-unit:cell}{cel} per sigaret, in zo'n $10^9$ \omterm{def:g10:cells-common-unit:cell}{cellen} van de
bekleding, over $\num{255000}$ sigaretten, is $10^{14}$ \omterm{def:g11:mutations:mutation}{mutaties} verspreid
over de long; enkele daarvan treffen p53 of een proto-oncogen in een \omterm{def:g10:cells-common-unit:cell}{cel}
die al een andere stap draagt. Het aantal pogingen maakt het
onwaarschijnlijke zeker.

\textbf{16.} Rokers: $\num{27000}$ sterfgevallen op $\num{11.25e6}$,
ongeveer 240 per 100\,000; niet-rokers: 3000 op $\num{33.75e6}$, ongeveer
9 per 100\,000.

\textbf{17.} Ongeveer 27 keer --- van dezelfde orde als de 17 tot 25 uit
de figuur voor zware rokers; het bevolkingscijfer mengt lichte en zware
rokers en ex-rokers.

\textbf{18.} Ongeveer 3000 per jaar, het aandeel van de niet-rokers. Niet
onmiddellijk: de \omterm{def:g11:mutations:mutation}{mutaties} die huidige en voormalige rokers al dragen,
zullen nog twintig of dertig jaar \omterm{def:g11:cancer:cancer}{kankers} opleveren.

\textbf{19.} Elke roker screenen bespaart ongeveer 20\% van
$\num{27000}$, zo'n 5000 sterfgevallen per jaar; stoppen met roken zou er
uiteindelijk ongeveer $\num{24000}$ voorkomen --- vijf keer zoveel.

\textbf{20.} De handtekening van G naar T bewees dat de \omterm{def:g11:mutations:mutation}{mutaties} van de
\omterm{def:g11:cancer:cancer}{tumor} uit tabaksrook kwamen; de \omterm{def:g11:cancer:cancer}{kanker} werd gemaakt door één geactiveerd
oncogen en het verlies van beide \omterm{def:g10:universal-dna:gene}{allelen} van p53; niet roken had haar
voorkomen.
\end{solution}
